% Moved to supplement from Results (operational case studies + compute profile).
\section{Operational Case Studies and Computational Profile}
\label{sup:operational}

\subsection{Operational Incident Evaluation: Event Metrics and Adaptation}
\label{subsec:incident_metrics}

Evaluating telemetry solely by point-level metrics can penalize slight temporal boundary offsets in continuous anomalies. In operational systems, faults represent extended continuous incidents; a detector that identifies the event with a minor phase lag is penalized on every out-of-boundary point.

To analyze operational incident detection, Table~\ref{tab:operational_incident_metrics} evaluates three complementary metrics across three telemetry streams: Room Occupancy (stream: \texttt{default}, 11 events), NASA SMAP (stream: \texttt{P-3}, 1 event), and CalIt2 (stream: \texttt{traffic}, 21 events). We compare strict Point-F1, Event-F1 ($P_{\text{event}}, R_{\text{event}}$), and Tatbul et al. Range-Aware F1 ($F1_{\text{range}}$) across three operational inference configurations: Standard (middle-window, static EVT), Sharp (dense causal, static EVT), and Adaptive (dense, dynamic rolling SPOT).

\begin{table*}[t]
\centering
\small
\caption{Operational incident detection comparison across representative telemetry streams: Room Occupancy (stream: \texttt{default}, 11 ground-truth events), NASA SMAP (stream: \texttt{P-3}, 1 thruster fault), and CalIt2 (stream: \texttt{traffic}, 21 pedestrian burst events). We contrast Standard (Middle, Static EVT), Sharp (Dense, Static EVT), and Adaptive (Dense, Dynamic SPOT) inference modes.}
\label{tab:operational_incident_metrics}
\resizebox{\linewidth}{!}{%
\begin{tabular}{@{}llccccccc@{}}
\toprule
\textbf{Evaluated Stream} & \textbf{Operational Configuration} & \textbf{Pt-F1} & \textbf{Pt-Prec} & \textbf{Evt-F1} & \textbf{Evt-Prec} & \textbf{Evt-Rec} & \textbf{Rng-F1} & \textbf{False Pos.} \\
\midrule
\multicolumn{9}{l}{\textit{Room Occupancy: Stream \texttt{default} (Environmental Sensing, 11 Ground-Truth Incidents)}} \\
\texttt{operator\_entropy} & Standard (Middle, Static EVT) & 0.3728 & 22.9\% & 0.5238 & 35.5\% & 100\% & 0.3708 & 3,842 \\
\texttt{operator\_entropy} & Sharp (Dense, Static EVT)   & 0.3329 & 20.0\% & \textbf{0.9167} & \textbf{84.6\%} & \textbf{100\%} & 0.3329 & 4,705 \\
\texttt{operator\_entropy} & Adaptive (Dense, Dynamic SPOT) & 0.1273 & 29.9\% & 0.6364 & 63.6\% & 63.6\% & 0.3699 & \textbf{223} ($\downarrow 21\times$) \\
\texttt{potential\_flow}   & Sharp (Dense, Static EVT)   & 0.3431 & 20.7\% & \textbf{0.9167} & \textbf{84.6\%} & \textbf{100\%} & 0.3431 & 4,496 \\
\texttt{potential\_flow}   & Adaptive (Dense, Dynamic SPOT) & 0.0441 & 10.2\% & 0.5714 & 60.0\% & 54.5\% & 0.1636 & \textbf{291} ($\downarrow 15\times$) \\
\midrule
\multicolumn{9}{l}{\textit{NASA SMAP: Stream \texttt{P-3} (Spacecraft Telemetry, 1 Continuous Fault Segment)}} \\
\texttt{operator\_entropy} & Standard (Middle, Static EVT) & 0.0087 & 14.6\% & 0.1818 & 10.0\% & 100\% & 0.2266 & 35 \\
\texttt{operator\_entropy} & Sharp (Dense, Static EVT)   & 0.1274 & 27.3\% & 0.1538 & 8.3\%  & 100\% & \textbf{0.3634} & 295 \\
\texttt{potential\_flow}   & Adaptive (Dense, Dynamic SPOT) & 0.0110 & 6.9\%  & \textbf{0.3333} & 20.0\% & \textbf{100\%} & 0.1213 & \textbf{108} \\
\midrule
\multicolumn{9}{l}{\textit{CalIt2: Stream \texttt{traffic} (Pedestrian Entrance Flow, 21 Transient Event Clusters)}} \\
\texttt{operator\_entropy} & Standard (Middle, Static EVT) & 0.0131 & 1.9\%  & 0.0588 & 7.7\%  & 4.8\%  & 0.0216 & 53 \\
\texttt{operator\_entropy} & Adaptive (Dense, Dynamic SPOT) & \textbf{0.1368} & 8.7\%  & \textbf{0.3226} & 50.0\% & 23.8\% & \textbf{0.1274} & 357 \\
\texttt{potential\_flow}   & Adaptive (Dense, Dynamic SPOT) & 0.0903 & 14.3\% & 0.0909 & \textbf{100.0\%} & 4.8\%  & 0.0714 & \textbf{42} \\
\bottomrule
\end{tabular}%
}
\end{table*}

Table~\ref{tab:operational_incident_metrics} highlights three operational characteristics:
First, static EVT thresholding achieves 100\% event recall on Room Occupancy and SMAP (\texttt{P-3}), but generates over 4,000 false alarm points on Room Occupancy due to unmodeled diurnal baseline shifts, causing point precision to collapse to $20.0\%$.
Second, dynamic rolling SPOT suppresses false alarms by $21\times$ (from 4,705 to 223 points) and raises precision to $29.9\%$, but event recall drops from 100\% to $63.6\%$. Tracking non-stationary drift via rolling quantiles creates an operational trade-off between alarm fatigue and missed events.
Third, on transient burst streams (CalIt2, 21 short spikes), models detect only $4.8\%\text{--}23.8\%$ of events under rolling SPOT, demonstrating that short pulses remain difficult to distinguish from ambient noise.


\subsection{Five-Stream Causal-vs-Buffered Operational Comparison}
\label{sup:causal_vs_buffered}

\label{subsec:causal_vs_buffered}

Table~\ref{tab:causal_vs_buffered} compares operational latency, false-alarm rates, and detection delay across five matched telemetry streams (NASA SMAP \texttt{P-3}, SMD \texttt{machine-1-2}, NASA MSL \texttt{M-1}, Daphnet \texttt{S01R01E1}, and GECCO \texttt{water\_quality}) under the strictly causal ($0$-lookahead) and buffered ($64$-lookahead) regimes defined in Section~\ref{subsec:causal_vs_buffered_protocol}.

\begin{table*}[t]
\centering
\small
\caption{Operational latency and detection fidelity under Strictly Causal ($0$-lookahead delay) versus Buffered Window ($64$-lookahead delay) inference across 5 representative telemetry streams. Under causal streaming, TS-JEPA achieves an empirical false-positive rate of $4.18\%$ ($\text{FAR}_{1\text{k}} = 41.76$ false alarms per 1k points), compared to $30.06\%$ for TimesNet and $20.81\%$ for TranAD.}
\label{tab:causal_vs_buffered}
\resizebox{\linewidth}{!}{%
\begin{tabular}{@{}llccccccc@{}}
\toprule
\textbf{Inference Setting} & \textbf{Model} & \textbf{Empirical FPR} & \textbf{$\text{FAR}_{1\text{k}}$ (/1k pts)} & \textbf{False Positives} & \textbf{Point-F1} & \textbf{Precision} & \textbf{Mean Delay ($\Delta t$)} & \textbf{Recall ($\tau \le 64$)} \\
\midrule
\multirow{3}{*}{\shortstack[l]{\textbf{Strictly Causal}\\($0$-Lookahead)}} 
& TS-JEPA & \textbf{0.0418} & \textbf{41.76} & \textbf{1,685.6} & 0.0688 & \textbf{0.1581} & 130.3 & 0.1071 \\
& TimesNet & 0.3006 & 300.63 & 16,952.0 & 0.0774 & 0.1525 & 251.2 & 0.2286 \\
& TranAD & 0.2081 & 208.13 & 16,637.6 & \textbf{0.0881} & 0.1255 & \textbf{43.2} & \textbf{0.3286} \\
\midrule
\multirow{3}{*}{\shortstack[l]{\textbf{Buffered Window}\\($64$-Lookahead)}} 
& TS-JEPA & \textbf{0.1514} & \textbf{151.39} & \textbf{2,696.0} & 0.0431 & 0.0355 & 89.6 & 0.1357 \\
& TimesNet & 0.2800 & 280.03 & 16,916.8 & \textbf{0.0807} & 0.1404 & 235.3 & 0.2571 \\
& TranAD & 0.1998 & 199.82 & 16,339.6 & 0.0197 & \textbf{0.2284} & \textbf{32.8} & \textbf{0.2786} \\
\bottomrule
\end{tabular}%
}
\end{table*}

Under strictly causal streaming, TS-JEPA maintains an empirical false-positive rate of $4.18\%$ ($41.76$ false alarms per $1{,}000$ points), generating an average of $1{,}685.6$ false alarms per stream. In contrast, TimesNet and TranAD generate $16{,}952.0$ and $16{,}637.6$ false alarms (empirical FPRs of $30.06\%$ and $20.81\%$). On GECCO water telemetry ($113{,}000$ test points), TimesNet and TranAD trigger over $81{,}500$ false alarms ($99.6\%$ FPR), whereas TS-JEPA triggers zero false alarms.

Two methodological trade-offs govern these results:
First, the false-alarm reduction under strict causality comes at the expense of early event capture. TS-JEPA detects $10.7\%$ of anomaly events within $\tau \le 64$ steps under causal streaming, compared to $22.9\%$ for TimesNet and $32.9\%$ for TranAD.

\subsection{Computational Profile: Complexity, Latency, and Memory Footprint}
\label{subsec:computational_profile}

Table~\ref{tab:computational_profile} benchmarks parameter complexity, peak GPU memory footprint (VRAM), forward inference latency, and streaming throughput across all evaluated architectures on an NVIDIA RTX 5090 GPU (input dimension $K=25$, batch size $32$, context $C=256$, suspect $S=64$).

\begin{table}[t]
\centering
\small
\caption{Computational profile, memory consumption, latency, and throughput benchmarked on an NVIDIA RTX 5090 GPU ($K=25$, window size $320$).}
\label{tab:computational_profile}
\resizebox{\linewidth}{!}{%
\begin{tabular}{@{}lccccc@{}}
\toprule
\textbf{Model Architecture} & \textbf{Parameters} & \textbf{VRAM (MB)} & \textbf{Latency (ms/win)} & \textbf{Infer (win/s)} & \textbf{Train (win/s)} \\
\midrule
TS-JEPA & 289,061 & \textbf{183.8} & 0.6072 & 1,646.8 & 612.8 \\
Operator-Entropy JEPA & \textbf{285,733} & 183.9 & 0.6939 & 1,441.0 & 553.2 \\
Reynolds-Stress JEPA & 295,953 & 184.0 & 0.7456 & 1,341.2 & 535.8 \\
Potential-Flow JEPA & 338,117 & 184.4 & 0.8445 & 1,184.2 & 400.2 \\
Helmholtz JEPA & 416,391 & 185.3 & 0.8378 & 1,193.6 & 526.6 \\
TimesNet & 648,553 & 279.6 & \textbf{0.2908} & 3,439.4 & 1,161.4 \\
TranAD & 205,746 & 1,280.0 & \textbf{0.1897} & \textbf{5,272.2} & \textbf{1,554.4} \\
\bottomrule
\end{tabular}%
}
\end{table}

As shown in Table~\ref{tab:computational_profile}, Operator-Entropy JEPA possesses fewer trainable parameters ($285{,}733$) than unconstrained TS-JEPA ($289{,}061$), because its explicit Koopman transition matrix $K \in \mathbb{R}^{32 \times 32}$ ($1{,}024$ weights) replaces the multi-layer MLP predictor. All JEPA variants maintain a compact peak GPU memory footprint of $183.8\text{--}185.3\,$MB, consuming nearly $7\times$ less memory than TranAD ($1{,}280.0\,$MB), whose self-attention maps scale quadratically with window length. Forward inference latency for the JEPA family ranges between $0.60\text{--}0.84\,$ms per window ($1{,}184\text{--}1{,}647$ windows per second). While TimesNet ($0.29\,$ms) and TranAD ($0.19\,$ms) achieve faster forward inference, the JEPA family operates well within sub-millisecond operational streaming deadlines on GPU hardware. These measurements characterize workstation accelerator throughput; on resource-constrained embedded processors or flight microcontrollers lacking dedicated tensor cores, sequential end-to-end latency depends directly on clock frequencies and memory bandwidth, motivating integer quantization (Section~\ref{sec:conclusion}) for edge deployment.
